\section{讨论}

\subsection{为什么做这项研究}

我们开展该研究的核心动机是：在 BRCA 亚型识别任务中，单组学模型虽可用，但对临床异质性的刻画能力有限。课程作业要求不仅是“训练一个分类器”，而是要证明多组学融合是否能在同一评估口径下提供更稳定、可解释的增益。因此，本研究重点放在“统一协议下的可复现比较”，而非单次切分下的偶然高分。

\subsection{方法差异的解释}

本轮结果显示 Concat 的均值最高，RNA-only 与 Stacking 次之，MOFA 波动最大。四种方法形成由浅入深的融合链条：RNA-only 提供单组学基线，Concat 引入跨模态互补，MOFA 在潜空间压缩冗余，Stacking 在决策层进一步融合。该排序说明“跨组学互补”已在当前数据上体现，但“复杂融合”不必然带来稳定增益。

\subsection{为什么严格 CV 是主结论口径}

在小样本高维任务中，单次切分结果方差较大，容易出现方法同分或排序反转。repeated CV 通过多次分层重采样提供更稳健的性能估计，因此更适合作为论文主证据。另一方面，若在同一数据上同时做模型选择与性能汇报，结论可能受到选择偏倚影响，这在方法学文献中已有系统讨论\cite{cawley2010}。因此，我们在实现层面把所有可学习步骤限定在训练折，并将测试折严格用于独立评估。

\subsection{MOFA 的价值与边界}

MOFA 的核心价值在于潜表示建模。我们已将 MOFA 纳入训练折拟合与测试折投影流程，保证评估一致性。但在当前样本规模下，MOFA 的方差较大（Accuracy std = 0.2193），说明其对因子数和折内样本扰动更敏感，需要通过更系统调参与更大样本来发挥优势。

从团队复盘角度，我们认为 MOFA 当前不占优并不等于“表示学习无效”，更可能反映了“模型复杂度与样本规模不匹配”。在 $n\ll p$ 的高维设置下，额外自由度会放大估计不稳定性\cite{hughes1968}，这也解释了为何简单且受约束的 Concat 在本数据上更稳。

\subsection{当前局限}

本文仍存在三点局限：

\begin{enumerate}
    \item 样本规模有限，HER2 类在当前对齐后被剔除，导致任务退化为三分类；
    \item 外部队列验证尚未完成；
    \item 超参数搜索空间仍可扩展。
\end{enumerate}

\subsection{后续工作}

下一步将沿三条路径推进：

\begin{enumerate}
    \item 扩展稳健性证据：增加随机种子、bootstrap 和外部验证，提升统计检验功效；
    \item 强化可解释性证据：围绕 LumB 易混类别开展错误归因、关键特征和模态贡献分析；
    \item 面向高分报告补强：补充失败案例讨论与生物学机制对照，形成“结果-机制-边界-改进”闭环。
\end{enumerate}

\subsection{本研究的 novelty 边界声明}

为避免过度主张，我们将本文 novelty 明确限定为“研究设计创新”而非“算法发明”：
\begin{enumerate}
    \item 在 BRCA 单癌种任务中实现了可审计的严格防泄露评估流程；
    \item 构建并实证比较了四层融合递进链条，给出可复算的性能趋势与统计边界；
    \item 将类别级误差（尤其 LumB 易错）纳入主讨论，形成面向后续改进的具体问题清单。
\end{enumerate}

\subsection{详细继续改进方案}

\begin{table}[H]
\centering
\caption{后续改进路线（高分导向）}
\label{tab:improvement-roadmap}
\begin{tabular}{p{0.17\linewidth}p{0.22\linewidth}p{0.24\linewidth}p{0.16\linewidth}p{0.14\linewidth}}
\toprule
阶段 & 目的 & 主要动作 & 预期收益 & 风险/代价 \\
\midrule
短期（仓库内可完成） & 提升证据密度 & 增加重采样轮次、补充 Balanced Accuracy 与 Macro-F1 显著性检验、完善类别级误差表 & 结论更稳健，报告说服力提升 & 计算时间增加 \\
中期（需额外数据） & 提升泛化可信度 & 引入外部 BRCA 队列或留出独立验证集，复现实验主线 & 缓解“仅单数据源有效”质疑 & 数据清洗与对齐成本高 \\
长期（研究深化） & 提升机制解释深度 & 联合通路分析/文献证据解释关键特征，探索图模型或更强表示学习 & 形成“性能+机制”双重贡献 & 实现复杂度与不确定性上升 \\
\bottomrule
\end{tabular}
\end{table}
